% Entire guide absent from English release/native project.
% Reconstructed from French original PDF pp.6--8, printed pp.1--3.
\chapter*{SGA 4、SGA \texorpdfstring{$4\frac12$}{4½} 与 SGA 5 的阅读线索}
\addcontentsline{toc}{chapter}{SGA 4、SGA 4½ 与 SGA 5 的阅读线索}
\label{reading-guide}
SGA~4 的第~I~至第~VI~讲给出 Grothendieck 拓扑的一般理论. 这些讲写得非常详细, 在研究某些特殊拓扑---例如产生晶体上同调的拓扑---时可能十分有用. 对于如此接近经典直觉的平展拓扑, 则不必设置这样庞大的护栏: 只须熟悉例如 Godement 的书 [4], 并怀有一点信心. 其他可用的参考文献还有 Artin 笔记 [1] 的第~I~至第~III~章、Giraud 的 Bourbaki 报告 [3], 或 [Arcata]~I. 第~VII~和第~VIII~讲开始研究平展拓扑, 比 Artin 的第~III~章和 [Arcata]~II 更为详细.

曲线的研究是平展上同调的关键. Artin 在 SGA~4 第~IX~讲开始这项研究; 第~X~讲给出关于上同调维数的推论. 其中的要点在 [Arcata]~III 中重新叙述. 第~XI~讲不是后文所必需的; 它在光滑复代数簇的特殊情形下, 给出了平展上同调与经典上同调之间比较定理的一个优美证明.

第~XII~和第~XIII~讲证明了基本的紧合态射换基定理. 正如 Artin 在 \emph{Algebraic approximation of structures over complete local rings}, \emph{Publ. Math. IHES} 36 (1969), 23--58, \S3 中指出的, 他的逼近定理可以简化这个证明. [Arcata]~IV 采用的就是这条路线; 其中只勾勒了证明所需的多次逐层约化, 对第~XIV~讲中的应用也只作了十分简短的提示. 这个定理使我们能够定义“紧支集上同调”, 以及“紧合支集高阶正像”函子 $\eR^i f_!$（这些函子的性质可归结为紧合态射的高阶正像函子的性质）. 这正是篇幅冗长的第~XVII~讲的主题; 要点已在 [Arcata]~IV 中说明, 但第~XVII~讲的一些部分也可能具有独立用途.

第~XV~和第~XVI~讲以光滑态射的基本局部无环性定理为中心. 在篇幅较简短的平行论述 [Arcata]~V 中, 可以看出 SGA~7 的影响; 该处也证明了所需的 SGA~2 中的结果.

[Arcata]~VI 专门讨论 Poincaré 对偶性; 它比 SGA~4 第~XVIII~讲更清楚, 但并不声称给出了完整证明. 对于坚持要求纯代数证明的读者, 通过引用 [Dualité]~\S3, 论证得以简化; 它没有证明迹态射所满足的形式体系（第~XVIII~讲~\S2）.

在平展上同调中, 存在类似于凝聚上同调的对偶形式体系. 为了建立它,\newline
Grothendieck 使用了奇点消解和纯性猜想（陈述见 [Cycle]~2.1.4）; 纯性猜想的相对情形在 SGA~4 第~XVI~讲中建立, 等特征情形则在假定奇点消解的条件下于 SGA~4 第~XIX~讲中建立. 对于维数为 $0$ 或 $1$ 的正则概形上的有限型概形, [Th. finitude] 用另一种方法建立了其中的关键结果. SGA~5 第~I~讲给出若干进一步的发展. SGA~5 第~III~讲说明怎样由这一形式体系推出十分一般的 Lefschetz--\allowbreak Verdier 迹公式.

由此可见, 在 SGA~5 的原始版本中, Lefschetz--\allowbreak Verdier 公式只是在猜想的条件下建立的. 此外, 其中没有计算局部项. 为了应用于 $L$ 函数, 这个讨论班给出了另一种证明; 在 Frobenius 态射的特殊情形下, 这个证明是完整的. [Rapport] 收录的就是这一证明. 其他参考文献包括: 关于陈述和化约框架, 参见 Grothendieck 的 Bourbaki 报告 [5]; 关于如何把关键情形化约到 Weil 已处理的一个情形, 参见 [2]~\S10 中对 Grothendieck 所作这一化约的简短说明; 关于后一情形的 $\ell$ 进处理, 参见 [Cycle]~\S3.

SGA~5 还详细讨论了从 $\dZ/\ell^n$ 系数上同调过渡到 $\dZ_\ell$ 系数上同调（由 J.~P. Jouanolou 撰写的第~V~和第~VI~讲）, 以及 Frobenius 态射的形式体系（由 C.~Houzel 撰写的第~XV~讲）. 由 I.~Bucur 撰写的第~X~和第~XII~讲, 计算曲线上一个层的 Euler--Poincaré 示性数, 以及某些对应在其上同调上定义的自同态的迹. Jouanolou 在第~VII~讲中计算经典概形的上同调并给出应用. 此外, Serre 的报告《Introduction à la théorie de Brauer》已收入他的书 [6] 的第三部分.

\section*{参考文献}
\begin{description}
\item[{[1]}] M.~Artin, \emph{Grothendieck topologies}, Harvard University, 1962.
\item[{[2]}] P.~Deligne, \emph{Les constantes des équations fonctionnelles des fonctions L}, 载 \emph{Modular functions of one variable II}, 501--597, Lecture Notes in Mathematics 349, Springer-Verlag.
\item[{[3]}] J.~Giraud, \emph{Analysis situs}, Séminaire Bourbaki 256, 1963 年 5 月, Benjamin.
\item[{[4]}] R.~Godement, \emph{Topologie algébrique et théorie des faisceaux}, Act. Sci. et Ind. 1252, Hermann, 1958.
\item[{[5]}] A.~Grothendieck, \emph{Formule de Lefschetz et rationalité des fonctions L}, Séminaire Bourbaki 279, 1964 年 12 月, Benjamin.
\item[{[6]}] J.~P. Serre, \emph{Représentations linéaires des groupes finis}, Collection Méthodes, Hermann, 1967.
\end{description}
\begin{description}
\item[SGA:] \emph{Séminaire de Géométrie algébrique du Bois-Marie}.
\item[SGA 4:] \emph{Théorie des topos et cohomologie étale des schémas},\\
由 M.~Artin、A.~Grothendieck 和 J.~L. Verdier 主持, Lecture Notes 269、270 和 305.
\item[SGA 5:] \emph{Cohomologie $\ell$-adique et fonctions L}, 由 A.~Grothendieck 主持, 即将在 Lecture Notes 出版.
\end{description}
